% Owner: Youssef. Sources: AGENTS.md (problem, idea), docs/threat-model.md (adversary, claim),
% docs/related-work.md section 3 (the narrowed contributions).
\section{Introduction}\label{sec:intro}

Coding agents now run shell commands on developers' machines. Given broad permissions, an agent
built on a large language model will install packages, delete build directories, rewrite
configuration and push to remote repositories, issuing each step as a line of shell. The
safeguard most agents ship is a list of forbidden command strings: block \cmd{rm -rf}, ask
before \cmd{sudo}. Such a list inspects the text of a command before the shell has read it, and
the shell is a language in which one effect can be written in many ways. The same deletion is
\cmd{rm -r dir}, \cmd{find dir -delete}, a Python one-liner, a base64 payload piped to \cmd{sh},
or \cmd{rm} spelled with quotes the shell removes and the filter does not. An industry study of
eleven open-source coding agents found that quoting and expansion tricks of this kind defeated
the command filters of ten of them, and that the one which held tokenised commands the way bash
itself does~\cite{guardfall}\checkcite{industry note, not peer reviewed}.

Evasion by a deliberate attacker is only half of the problem. The other half is that a refusal
does not end the interaction. An agent has a task, the obvious command was refused, and it
reaches for the next way to finish the job: \cmd{rm} is refused, so it tries \cmd{find -delete}.
A recent survey of agent safety calls the cost of this the \emph{verifier tax}, reporting one
study in which blocking 94\% of unsafe actions still left under 5\% of tasks completed both
safely and correctly, because the agents found other unsafe routes~\cite{survey-safe-agents}%
\checkcite{abstract only; the survey rates this evidence as very low quality}. A guard that
judges each command in isolation cannot see this pattern at all, because every command in it
looks new.

We treat the problem as one of compilation. \sys{} sits between the agent and the shell and
compiles each command line, as a program in a precisely specified subset of bash, into the set of
effects it would have: the files it would read, write or delete, the hosts it would contact, the
programs it would hand code to, the privileges it would change. A policy is then written over
those effects and the regions of the filesystem they fall in, not over strings, so the many
spellings of one deletion meet the same rule. The compiler has the usual phases. A lexer records,
for every piece of every word, how it was quoted and whether its value is known from the source
alone. A recursive-descent parser builds a syntax tree for the accepted language, parsing command
substitutions recursively rather than treating them as opaque. A resolver and an effect analysis
turn words into paths, using a session state that remembers the working directory and the
variables earlier commands set. And an emitter, the code-generation phase, writes the analysed
tree back out as a canonical command line, which is what the gate executes instead of the agent's
original text. Whatever our reading of a command, the shell runs the command we read.

Around the compiler sit two further components. A zero-shot classifier, built on an open-weight
natural language inference model, compares each command that follows a refusal with the command
that was refused and judges whether it pursues the same effect by other means. It may only
tighten a verdict, never loosen one. And a gate runs the whole pipeline for each command, tells
the agent why a command was refused, and executes allowed commands inside a disposable sandbox
that also records what they changed.

The research question is whether an effect-based compiler with a retry-aware classifier stops
agent workarounds better than a string denylist, at an acceptable rate of false refusals on
benign developer tasks. Our primary adversary is the over-eager agent described above, one that
pursues its task past a refusal without concealing what it does. Against an agent following
injected instructions, only the compiler's structural analysis holds, and we do not claim
otherwise for the classifier.

This paper makes three contributions.
\begin{enumerate}
  \item \textbf{A specified command language with quoting provenance and canonical emission.}
    We publish the subset of bash we accept, as a grammar, and refuse everything outside it
    with a named diagnostic. Every piece of every word carries how it was written, so a value
    that depends on the environment is marked unknown rather than guessed, and the command that
    runs is our re-serialisation of the tree we checked. To our knowledge no agent command guard
    states the language it accepts, and none executes its own output.
  \item \textbf{An effect-and-zone policy with bounded path sets and session state.} Targets
    are resolved to an exact path, a directory bounding every path touched, or unknown, which
    separates \cmd{rm -rf build} inside the workspace from \cmd{find / -delete} without refusing
    both. Paths are resolved against the session's working directory, so a \cmd{cd} out of the
    workspace followed by a relative read is not invisible, as it is to a guard that judges
    each command alone.
  \item \textbf{A study of agent behaviour after a refusal}, measuring how often agents work
    around a refused effect, whether the history of refusals helps detect it, and whether what
    the refusal message says changes what the agent does next.
\end{enumerate}

The guarantee we claim is deliberately narrow. For every command the gate allows, the direct
effects of the pinned version of bash evaluating that command line are within policy, provided
the entries of our table of command effects are correct and the filesystem does not change
between the check and the execution. We do not claim to follow effects through programs the
agent writes and later runs, to cover agent tools other than the shell, or to protect the
workspace from the agent it was given to.

\Cref{sec:related} places this work among command guards for agents and the formal study of the
shell. \Cref{sec:design} presents the design, \cref{sec:evaluation} the evaluation, and
\cref{sec:discussion} its limits.
